\section{Discussão}
\label{sec:discussao}

\textbf{Camuflagem.} É o principal erro que sobrou. Onde a camiseta
escura passa na frente do ventilador preto, ou o short na frente da
cadeira, as cores ficam a uma distância $D^2$ menor que o limiar, e o
MOG2 classifica aquele pedaço da pessoa como fundo
(Seção~\ref{sec:mog2}). Como o método só compara cores, nenhum
parâmetro resolve isso sem criar outro erro. Uma melhoria seria usar
também textura ou bordas, que mudam quando a pessoa entra mesmo com cor
parecida. Na demonstração, a solução prática é tirar objetos escuros de
trás da pessoa.

\textbf{Todo ajuste é uma troca.} Com o \texttt{varThreshold} em 16, o
MOG2 fica mais sensível e tapa parte da camuflagem: os acertos nos
vídeos subiram de 4 para 8 (IoU médio 0,821). Na gravação da webcam,
porém, os quadros com mais de um objeto passaram de 14 para 128, em
603: a pessoa perto da câmera muda levemente a luz da cama ao lado
(sombra e reflexo), e com o limiar baixo isso já vira mancha. Da mesma
forma, sem a etapa 6 os acertos subiriam de 4 para 6, porque um buraco
de camuflagem dentro da silhueta é igual, para o programa, ao vão entre
o braço e o corpo; mas a parede antiga voltaria a aparecer nesse vão.
Mantivemos os valores testados ao vivo. Escolher o limiar olhando os
mesmos 12 quadros da avaliação, aliás, deixaria o número otimista.

\textbf{Contagem.} Os três erros de contagem foram o braço levantado
separado do corpo, quase sempre porque a camiseta entre eles sumiu por
camuflagem. Juntar contornos próximos resolveria esses casos, mas
juntaria também duas pessoas lado a lado.

\textbf{O gabarito também erra.} Na primeira avaliação, o IoU médio foi
0,769. A sobreposição mostrou que o gabarito passava de 4 a 8\,px para
fora da pessoa: a ferramenta mostrava o quadro pequeno (269\,px de
largura) e o pincel era grande. Com zoom e um modo de correção na
ferramenta, o grupo corrigiu as bordas, e o IoU médio foi para 0,791,
sem nenhuma mudança no programa.

\textbf{Mudança de luz.} Congelar o modelo faz uma pessoa parada
continuar aparecendo, mas impede o MOG2 de se adaptar se a luz do
ambiente mudar. Uma melhoria seria atualizar o modelo devagar, só fora
da máscara da pessoa; hoje, a tecla B reaprende o fundo. Pelo mesmo
motivo, a cena precisa estar vazia no aprendizado: nos primeiros testes
com a webcam, a pessoa ficou na frente da câmera, virou parte do fundo
e só os contornos dela apareciam, o que motivou a contagem regressiva.
